\chapter{Extremal Problems and Ramsey Theory}

\begin{center}\itshape
Complete disorder is impossible.\\
\upshape(Theodore Motzkin)
\end{center}

In December 1928 a young Cambridge philosopher, Frank Ramsey, read a paper to the London Mathematical Society called
\emph{On a problem of formal logic}. He wanted to decide which sentences of a certain logical form are true, and on the way
he needed a small combinatorial lemma: if you colour the pairs of an infinite set with two colours, some infinite subset has
all its pairs of the same colour. The logic problem was later shown to be hopeless in general, by Church and Turing. The lemma
became a field. Ramsey never saw it: he died in January 1930, at the age of twenty-six, having also founded the subjective theory of
probability and written two papers on economics that are still taught.

Ramsey's lemma says that \emph{every} large structure, however it is coloured, contains a perfectly ordered piece. A dozen years later,
Paul Tur\'an, then held in a Hungarian forced labour camp, asked a cousin of the same question: how many edges can a graph
have before it is forced to contain a complete subgraph $K_r$? He found the answer there, with no books and little paper, and published it in 1941. Tur\'an's question starts \emph{extremal graph theory}; Ramsey's starts \emph{Ramsey
theory}. Both ask how much disorder a finite world can hold.

The answers are often surprisingly exact for small cases and surprisingly unknown for slightly larger ones. Tome~3, Chapter~12
proved that among any six people there are three mutual friends or three mutual strangers, and that the number $R(k)$ needed for $k$
instead of three lies between $2^{k/2}$ and $4^k$. Paul Erd\H os liked to tell the following story. Suppose aliens invade the Earth and
threaten to destroy it within a year unless we tell them $R(5)$. We should put all our mathematicians and computers to work, and we would
probably succeed. If they asked for $R(6)$, we should try to destroy the aliens.

This chapter proves Tur\'an's theorem twice, finds the extremal number of the square $C_4$ with a projective plane, computes
$R(3,4)$, $R(3,5)$ and $R(4,4)$ with three beautiful colourings, proves Ramsey's theorem for sets of any size by infinity and
compactness, and then moves to the integers: Schur's theorem, which shows that Fermat's last theorem cannot be proved modulo
primes, van der Waerden's theorem on arithmetic progressions, and the sunflower lemma, which the theory of computation will use to
prove that some functions need huge circuits.

\begin{think}
Colour the numbers $1,2,\dots,9$ red and blue. Can you avoid three numbers of the same colour in arithmetic progression, like $2,5,8$?
What about $1,\dots,8$?
\end{think}

\section{Forbidding a triangle, forbidding a clique}

For a graph $H$, the \emph{extremal number} $\mathrm{ex}(n,H)$ is the largest number of edges of a graph on $n$ vertices that contains no
copy of $H$ as a subgraph. The first case is the triangle.

\begin{theorem}[Mantel, 1907]\label{thm:mantel}
A graph with $n$ vertices and no triangle has at most $\frac{n^2}4$ edges.
\end{theorem}
\begin{proof}
Let $G$ have $m$ edges and no triangle. For an edge $xy$, the neighbourhoods $N(x)$ and $N(y)$ are disjoint, since a common neighbour would close a
triangle; so $d(x)+d(y)\le n$. Sum over the $m$ edges: a vertex $v$ appears in $d(v)$ of them, so the left side becomes $\sum_vd(v)^2$, and
$\sum_vd(v)^2\le nm$. By the Cauchy--Schwarz inequality, $\sum_vd(v)^2\ge\frac1n\bigl(\sum_vd(v)\bigr)^2=\frac{4m^2}n$. Hence $\frac{4m^2}n\le nm$ and
$m\le\frac{n^2}4$.
\end{proof}

The complete bipartite graph $K_{\lfloor n/2\rfloor,\lceil n/2\rceil}$ has no triangle (a triangle is an odd cycle, Theorem~\ref{thm:bipartite}) and
$\lfloor\frac{n^2}4\rfloor$ edges, so $\mathrm{ex}(n,K_3)=\lfloor\frac{n^2}4\rfloor$. To forbid $K_{r+1}$, the natural guess is to split the vertices into $r$
classes and join every pair of vertices in different classes.

\begin{definition}\label{def:turan}
The \emph{Tur\'an graph} $T(n,r)$ is the complete $r$-partite graph on $n$ vertices whose classes have sizes $\lfloor\frac nr\rfloor$ or
$\lceil\frac nr\rceil$. Its number of edges is denoted $t(n,r)$.
\end{definition}

A complete multipartite graph whose classes have sizes $n_1,\dots,n_r$ (some possibly $0$) has $\frac12(n^2-\sum_in_i^2)$ edges: from all ordered pairs of
vertices, remove those in the same class, and halve. A clique has at most one vertex in each class, so $T(n,r)$ contains no $K_{r+1}$.

\begin{lemma}\label{lem:balance}
Among complete multipartite graphs on $n$ vertices with at most $r$ classes, $T(n,r)$ is, up to isomorphism, the only one with the most edges.
\end{lemma}
\begin{proof}
Use $r$ classes, allowing empty ones. If two class sizes $a\ge b+2$ exist, move one vertex from the class of size $a$ to the class of size $b$:
the moved vertex loses $b$ neighbours and gains $a-1$, a net gain of $a-1-b\ge1$. So a graph with the most edges has all class sizes within $1$ of each other.
Sizes $n_1,\dots,n_r$ with $|n_i-n_j|\le1$ and sum $n$ are determined up to order: if $n=qr+s$ with $0\le s<r$, exactly $s$ of them equal $q+1$ and the others $q$.
\end{proof}

\begin{figure}[ht]
\centering
\fig{ch08-turan}
\caption{The Tur\'an graph $T(10,3)$: classes of sizes $4$, $3$, $3$ (blue, red, green), every pair of vertices in different classes joined. It has $33$
edges and no $K_4$, and every graph on $10$ vertices with $34$ edges contains a $K_4$.}
\label{fig:turan}
\end{figure}

\begin{theorem}[Tur\'an, 1941]\label{thm:turan}
Let $r\ge1$. A graph with $n$ vertices and no $K_{r+1}$ has at most $t(n,r)$ edges, with equality only for $T(n,r)$.
\end{theorem}
\begin{proof}
The proof, due to Zykov (1949), never looks at the graph as a whole: it improves it one vertex at a time, by making a vertex a twin of another.
Let $G$ have no $K_{r+1}$ and the largest possible number $e(G)$ of edges among such graphs on $n$ vertices. We show that non-adjacency is an
equivalence relation on its vertices: being non-adjacent is reflexive (no loops) and symmetric, so the claim is transitivity.

\emph{Cloning.} For a vertex $v$ and another vertex $u$ not adjacent to $v$, let $G_{u\to v}$ be the graph obtained by deleting all edges at $u$ and joining $u$ to
every vertex of $N(v)$; then $u$ and $v$ are non-adjacent twins. A clique of $G_{u\to v}$ contains at most one of $u,v$, and replacing $u$ by $v$ turns
it into a clique of $G$ of the same size; so $G_{u\to v}$ has no $K_{r+1}$. Its number of edges is $e(G)-d(u)+d(v)$.

Suppose $u\not\sim v$, $v\not\sim w$ but $u\sim w$. If $d(u)>d(v)$, then $G_{v\to u}$ has more edges than $G$, which is impossible; so $d(u)\le d(v)$, and likewise
$d(w)\le d(v)$. Now clone $v$ twice: delete all edges at $u$ and at $w$ (that is $d(u)+d(w)-1$ edges, as $uw$ is counted twice) and join both $u$ and $w$ to
$N(v)$. In the new graph $u,v,w$ are pairwise non-adjacent twins, so again a clique contains at most one of them and there is no $K_{r+1}$. The number of
edges is $e(G)-d(u)-d(w)+1+2d(v)\ge e(G)+1$, a contradiction.

So the classes of non-adjacency partition the vertices, and $G$ is complete multipartite: two vertices are adjacent exactly when they are in different classes.
Taking one vertex per class gives a clique, so there are at most $r$ classes, and by Lemma~\ref{lem:balance}, $e(G)\le t(n,r)$, with equality only for
$T(n,r)$. As $T(n,r)$ has no $K_{r+1}$, the maximum is $t(n,r)$. Finally, any $K_{r+1}$-free graph with $t(n,r)$ edges is itself one of the graphs with the most edges,
so the argument applies to it, and it is $T(n,r)$.
\end{proof}

Here is a second proof of the bound, by probability, which loses only the rounding. Tome~3, Exercise~12.1 shows (by keeping the vertices that come
before all their neighbours in a random order) that every graph $H$ has an independent set of size at least $\sum_v\frac1{d_H(v)+1}$, the
\emph{Caro--Wei bound}.

\begin{proposition}\label{prop:turanprob}
A graph with $n$ vertices, $m$ edges and no $K_{r+1}$ satisfies $m\le\bigl(1-\frac1r\bigr)\frac{n^2}2$.
\end{proposition}
\begin{proof}
Apply Caro--Wei to the complement $H=\bar G$: its independent sets are the cliques of $G$, of size at most $r$, and $d_H(v)+1=n-d_G(v)$. So
\[
r\ge\sum_v\frac1{n-d_G(v)}\ge\frac{n^2}{\sum_v(n-d_G(v))}=\frac{n^2}{n^2-2m},
\]
by the Cauchy--Schwarz inequality $\bigl(\sum_v1\bigr)^2\le\sum_v\frac1{x_v}\sum_vx_v$ for positive $x_v$. Hence $n^2-2m\ge\frac{n^2}r$.
\end{proof}

When $r$ divides $n$, the bound is exactly $t(n,r)$. With $r=2$ it is Mantel's theorem again.

\begin{remark}
For a general graph $H$ with chromatic number $\chi(H)=r+1\ge3$, the Tur\'an graph $T(n,r)$ does not contain $H$ (its vertices would be coloured by the
$r$ classes), so $\mathrm{ex}(n,H)\ge t(n,r)$. The Erd\H os--Stone theorem (1946), in the form given by Erd\H os and Simonovits, says that this is the truth
up to lower order:
\[
\mathrm{ex}(n,H)=\Bigl(1-\frac1{\chi(H)-1}+o(1)\Bigr)\frac{n^2}2.
\]
The chromatic number alone decides the density. For bipartite $H$ the theorem only says $\mathrm{ex}(n,H)=o(n^2)$, and the true order of magnitude is
known only for a few bipartite graphs: the subject of the next section.
\end{remark}

\section{Forbidding a square}

How many edges can a graph on $n$ vertices have without containing a square $C_4$? Equivalently, how many friendships can there be among $n$
people if no two people have two common friends? The answer is of order $n^{3/2}$: the upper bound is a counting argument, and the lower bound comes
from geometry over a finite field.

\begin{theorem}[K\H ov\'ari, S\'os and Tur\'an, 1954]\label{thm:kst}
Let $1\le s\le t$. A graph with $n$ vertices and no subgraph $K_{s,t}$ ($s$ vertices all joined to the same $t$ vertices) has at most
\[
\tfrac12(t-1)^{1/s}n^{2-1/s}+\tfrac12(s-1)n
\]
edges. In particular $\mathrm{ex}(n,C_4)\le\frac12n^{3/2}+\frac12n$.
\end{theorem}
\begin{proof}
Count the pairs $(v,S)$ where $v$ is a vertex and $S$ is a set of $s$ neighbours of $v$. For each $v$ there are $\binom{d(v)}s$ of them. For each $s$-set $S$, the
vertices $v$ adjacent to all of $S$ number at most $t-1$, since $t$ of them together with $S$ would form a $K_{s,t}$. So
\[
\sum_v\binom{d(v)}s\le(t-1)\binom ns\le(t-1)\frac{n^s}{s!}.
\]
Let $f(x)=\frac1{s!}x(x-1)\cdots(x-s+1)$ for $x\ge s-1$ and $f(x)=0$ for $x<s-1$; then $f(d)=\binom ds$ for every integer $d\ge0$. On $[s-1,\infty)$, $f$ is a product of
$s$ non-negative, non-decreasing affine functions, so it is convex (the second derivative of a product of such functions is a sum of products of
non-negative terms) and non-decreasing; glued to $0$ on the left, with right derivative $f'(s-1)\ge0$, it is convex on $\R$. With $\bar d=\frac{2m}n$ the
average degree, Jensen's inequality gives $nf(\bar d)\le\sum_vf(d(v))\le(t-1)\frac{n^s}{s!}$. If $\bar d<s-1$ then $m<\frac12(s-1)n$ and we are done. Otherwise
$f(\bar d)\ge\frac{(\bar d-s+1)^s}{s!}$, so $n(\bar d-s+1)^s\le(t-1)n^s$, that is $\bar d\le(t-1)^{1/s}n^{1-1/s}+s-1$. Multiply by $\frac n2$. A square is $K_{2,2}$.
\end{proof}

For the lower bound we need many edges and no square, that is, any two vertices with at most one common neighbour. Two points of a plane
lie on one line: this is the source of the construction. Let $p$ be a prime and $\mathbb F_p=\Z/p\Z$. The \emph{projective plane} over $\mathbb F_p$ has as
points the lines through the origin in $\mathbb F_p^3$, that is, the non-zero vectors up to a non-zero scalar factor. There are $\frac{p^3-1}{p-1}=p^2+p+1$
of them. (Chapter~12 studies these planes for their own sake.)

\begin{theorem}[Erd\H os, R\'enyi and S\'os, 1966; Brown, 1966]\label{thm:polarity}
Let $p$ be a prime and $n=p^2+p+1$. The graph $G_p$ whose vertices are the points of the projective plane over $\mathbb F_p$, with $x\sim y$ when $x\ne y$ and
$x\cdot y=x_1y_1+x_2y_2+x_3y_3=0$, contains no $C_4$ and has at least $\frac12p(p^2+p+1)$ edges. Hence $\mathrm{ex}(n,C_4)\ge\frac12(1-o(1))n^{3/2}$.
\end{theorem}
\begin{proof}
Adjacency is well defined, since $x\cdot y=0$ does not depend on the scalar multiples chosen. A square $x\,z\,y\,z'$ would give two distinct points
$x\neq y$ with two common neighbours $z\ne z'$. Representatives of two distinct points are linearly independent, so the $2\times3$ matrix with rows $x$
and $y$ has rank $2$ and its kernel, the set of $z$ with $z\cdot x=z\cdot y=0$, has dimension $1$: it is a single point. So two points have at most one common
neighbour, and $G_p$ has no $C_4$.

The neighbours of $x$ are the points of the plane $x^\perp=\{z:z\cdot x=0\}$, a $2$-dimensional subspace containing $\frac{p^2-1}{p-1}=p+1$ points, minus $x$
itself if $x\cdot x=0$. So every degree is at least $p$, and there are at least $\frac12p(p^2+p+1)$ edges. Since $n^{3/2}=(p^2+p+1)^{3/2}=p^3(1+O(\frac1p))$, this is
$\frac12(1-o(1))n^{3/2}$ along these values of $n$; for other $n$, use the largest $p$ with $p^2+p+1\le n$ and the fact that consecutive primes $p<p'$ satisfy
$p'/p\to1$ (the prime number theorem, Tome~5).
\end{proof}

So $\mathrm{ex}(n,C_4)\sim\frac12n^{3/2}$: the counting bound and the geometric construction agree to the constant. The same idea with spheres in
three-dimensional space over a finite field gives $\mathrm{ex}(n,K_{3,3})\asymp n^{5/3}$ (Brown, 1966). Beyond that, much is: whether $\mathrm{ex}(n,K_{4,4})$ has order $n^{7/4}$, as Theorem~\ref{thm:kst} suggests, or $\mathrm{ex}(n,C_8)$ order $n^{5/4}$, nobody knows. The bipartite
case is the frontier of the subject.

\section{Ramsey numbers}

Write $R(k,l)$ for the smallest $n$ such that every colouring of the edges of $K_n$ in red and blue contains a red $K_k$ or a blue $K_l$. Tome~3, Proposition~12.1
proved the recursion $R(k,l)\le R(k-1,l)+R(k,l-1)$ (look at the red and blue neighbours of one vertex), hence $R(k,l)\le\binom{k+l-2}{k-1}$, and $R(3,3)=6$; Theorem~12.2 there
proved $R(k,k)>2^{k/2}$ with a random colouring. Clearly $R(k,l)=R(l,k)$, $R(2,l)=l$, and $R(k,l)>n$ means that there is a \emph{good} colouring of $K_n$, with no red
$K_k$ and no blue $K_l$. One extra observation, about parity, already gives exact values.

\begin{proposition}\label{prop:ramseyparity}
If $a=R(k-1,l)$ and $b=R(k,l-1)$ are both even, then $R(k,l)\le a+b-1$.
\end{proposition}
\begin{proof}
Let $n=a+b-1$ and suppose a colouring of $K_n$ has no red $K_k$ and no blue $K_l$. A vertex $v$ has at most $a-1$ red neighbours: $a$ of them would contain a
blue $K_l$, or a red $K_{k-1}$ which with $v$ would make a red $K_k$. Likewise it has at most $b-1$ blue neighbours. As $v$ has $n-1=(a-1)+(b-1)$ neighbours, it has
exactly $a-1$ red ones. Then the red graph has $n$ vertices of odd degree $a-1$, and $n=a+b-1$ is odd: the sum of degrees $n(a-1)$ is odd, while it is twice the number
of red edges. This contradiction proves the claim.
\end{proof}

The circulant graph on $\Z/n\Z$ with \emph{connection set} $S=-S\subseteq\Z/n\Z\setminus\{0\}$ joins $x$ and $y$ when $x-y\in S$. Colour its edges red and the other edges of $K_n$ blue.
Translations $x\mapsto x+c$ preserve both colours, so a monochromatic clique can always be moved to contain $0$: this halves the work in every check below.

\begin{theorem}[Greenwood and Gleason, 1955]\label{thm:smallramsey}
$R(3,4)=9$, $R(3,5)=14$ and $R(4,4)=18$.
\end{theorem}
\begin{proof}
\emph{Upper bounds.} $R(3,3)=6$ and $R(2,4)=4$ are even, so $R(3,4)\le6+4-1=9$ by Proposition~\ref{prop:ramseyparity}. Then $R(3,5)\le R(2,5)+R(3,4)\le5+9=14$, and
$R(4,4)\le R(3,4)+R(4,3)\le18$.

\emph{$R(3,4)>8$.} On $\Z/8\Z$ take $S=\{\pm1,\pm4\}$ (Figure~\ref{fig:circulants}, left). A red triangle containing $0$ would be $\{0,a,b\}$ with $a,b,b-a\in S=\{1,4,7\}$; the
differences of two elements of $\{1,4,7\}$ are $\pm3$ and $\pm6$, never in $S$. A blue $K_4$ containing $0$ would consist of $0$ and three elements of $\{2,3,5,6\}$ with
pairwise differences in $\{2,3,5,6\}$. Among $2,3,5,6$, the pairs with such a difference are $\{2,5\}$, $\{3,5\}$, $\{3,6\}$, and no three elements are pairwise of this kind.

\emph{$R(3,5)>13$.} On $\Z/13\Z$ take $S=\{\pm1,\pm5\}=\{1,5,8,12\}$ (Figure~\ref{fig:circulants}, middle). The sums $a+b$ of two elements of $S$ are $2,6,9,0,10,0,4,3,7,11$, never in
$S$, so there is no red triangle $\{0,a,a+b\}$. A blue $K_5$ containing $0$ consists of $0$ and four pairwise blue vertices of $B=\{2,3,4,6,7,9,10,11\}$. The red edges
inside $B$ are
\[
2\,3,\;3\,4,\;4\,9,\;9\,10,\;10\,11,\;11\,6,\;6\,7,\;7\,2\quad\text{and}\quad2\,10,\;3\,11.
\]
The first eight form a cycle of length $8$ through all of $B$, and the only two sets of four pairwise non-adjacent vertices on an $8$-cycle are the alternate
ones, $\{2,4,10,6\}$ and $\{3,9,11,7\}$. The first contains the red edge $2\,10$, the second $3\,11$. So there is no blue $K_5$.

\emph{$R(4,4)>17$.} The squares modulo $17$ are $Q=\{1,2,4,8,9,13,15,16\}$. Since $-1=16$ is a square, $Q=-Q$, and the \emph{Paley graph} $P_{17}$ is the circulant graph
with $S=Q$ (Figure~\ref{fig:circulants}, right). Multiplication by a non-zero square $c$ is a bijection of $\Z/17\Z$ that maps $Q$ to $Q$, so it preserves the colours. A red $K_4$ can be
translated to contain $0$, then its other vertices are squares, and multiplying by the inverse of one of them (a square) we may assume it contains $0$ and $1$. Its
two other vertices are then in $Q\cap(Q+1)$, adjacent to both $0$ and $1$. Now $Q+1=\{2,3,5,9,10,14,16,0\}$, so $Q\cap(Q+1)=\{2,9,16\}$, and the differences $7,14,7$ of these
three numbers are not squares: no two of them are red-adjacent, and there is no red $K_4$. Finally multiplication by the non-square $3$ maps $Q$ onto the non-squares, so it
maps red edges to blue edges and blue to red: the blue graph is isomorphic to the red one, and has no $K_4$ either.
\end{proof}

\begin{figure}[ht]
\centering
\fig{ch08-circulants}
\caption{The red graphs of three extremal colourings: on $\Z/8\Z$ with differences $\pm1,\pm4$; on $\Z/13\Z$ with differences $\pm1,\pm5$; on $\Z/17\Z$ with the square
differences (the Paley graph $P_{17}$). All missing edges are blue. They show $R(3,4)>8$, $R(3,5)>13$, $R(4,4)>17$.}
\label{fig:circulants}
\end{figure}

\begin{remark}
These numbers fit a pattern: the colourings are \emph{algebraic}. The differences $\pm1,\pm5$ modulo $13$ are exactly the non-zero cubes (Exercise~\ref{ex:fermatp}), and the
square differences give $P_{17}$. Beyond them, algebra gives out. The values $R(3,6)=18$, $R(3,7)=23$, $R(3,8)=28$, $R(3,9)=36$ and $R(4,5)=25$ needed long computer
searches, the last one (McKay and Radziszowski, 1995) years of them; Radziszowski's survey \emph{Small Ramsey numbers} keeps the list. For $R(5,5)$ the best
known bounds are $43\le R(5,5)\le46$, the upper one from 2024 (Angeltveit and McKay). The aliens have not yet asked.
\end{remark}

With more colours, the same neighbourhood argument works. Write $R_r(3)$ for the smallest $n$ such that every colouring of the edges of $K_n$ with $r$ colours has a monochromatic
triangle.

\begin{proposition}\label{prop:r333}
$R_3(3)\le17$, and in general $R_r(3)\le r\bigl(R_{r-1}(3)-1\bigr)+2$.
\end{proposition}
\begin{proof}
Let $n=r(R_{r-1}(3)-1)+2$ and colour $K_n$ with $r$ colours. A vertex $v$ has $n-1=r(R_{r-1}(3)-1)+1$ edges, so by the pigeonhole principle at least $R_{r-1}(3)$ of them
have the same colour, say colour $r$. If two of the corresponding neighbours are joined in colour $r$, they form a triangle of colour $r$ with $v$. Otherwise the
edges among these $R_{r-1}(3)$ neighbours use only $r-1$ colours, and contain a monochromatic triangle. For $r=3$: $3\cdot5+2=17$.
\end{proof}

A colouring of $K_{16}$ with three colours and no monochromatic triangle exists (Exercise~\ref{ex:r333lower}), so $R_3(3)=17$. For four colours, only
$51\le R_4(3)\le62$ is known.

\section{Ramsey's theorem}

Colour pairs, or triples, or $k$-element sets. The \emph{$k$-subsets} of a set $X$ form the set $\binom Xk$. A colouring $c:\binom Xk\to\{1,\dots,r\}$ is
\emph{constant} on $Y\subseteq X$ if all $k$-subsets of $Y$ get the same colour; then $Y$ is \emph{homogeneous}. Ramsey's original statement is infinite, and
infinite is easier.

\begin{theorem}[Ramsey, 1930]\label{thm:ramseyinf}
Let $k,r\ge1$ and let $X$ be an infinite set. For every colouring of $\binom Xk$ with $r$ colours there is an infinite homogeneous subset $Y\subseteq X$.
\end{theorem}
\begin{proof}
Induction on $k$. For $k=1$ this is the infinite pigeonhole principle: one of the finitely many colour classes of the points is infinite.

Let $k\ge1$, assume the theorem for $k$, and let $c$ colour the $(k+1)$-subsets of $X$. We build points $x_1,x_2,\ldots$ and infinite sets $X=X_0\supseteq X_1\supseteq\cdots$ as
follows. Given the infinite set $X_{i-1}$, pick any $x_i\in X_{i-1}$, and colour each $k$-subset $S$ of $X_{i-1}\setminus\{x_i\}$ with $c(S\cup\{x_i\})$. By the induction hypothesis,
there is an infinite subset $X_i\subseteq X_{i-1}\setminus\{x_i\}$ on which this colouring is constant, say with colour $\chi_i$. The points $x_1,x_2,\ldots$ are distinct, and
$x_j\in X_i$ for $j>i$.

Now take indices $i<j_1<\dots<j_k$. The points $x_{j_1},\dots,x_{j_k}$ lie in $X_i$, so $c(\{x_i,x_{j_1},\dots,x_{j_k}\})=\chi_i$: the colour of a $(k+1)$-subset of $\{x_1,x_2,\ldots\}$
depends only on its element of smallest index. By the pigeonhole principle, some colour $\chi$ equals $\chi_i$ for infinitely many $i$; the set $Y$ of these $x_i$ is
infinite and homogeneous of colour $\chi$.
\end{proof}

Chapter~1 promised a deduction of the finite theorem from the infinite one. Here it is.

\begin{theorem}[finite Ramsey theorem]\label{thm:ramseyfin}
For all $k,r,m\ge1$ there is an $N$ such that every colouring of the $k$-subsets of $\{1,\dots,N\}$ with $r$ colours has a homogeneous set of size $m$. The
least such $N$ is denoted $R^{(k)}_r(m)$.
\end{theorem}
\begin{proof}
Suppose not: for each $N$ there is a \emph{bad} colouring of $\binom{\{1,\dots,N\}}k$, with no homogeneous $m$-set. Take a propositional variable $x_{S,j}$ for every $k$-subset $S$
of $\N^*=\{1,2,\ldots\}$ and every colour $j$; there are countably many. Let $\Gamma$ contain, for each $S$, the formula $x_{S,1}\lor\dots\lor x_{S,r}$, and for each $m$-subset $M$ of $\N^*$ and each colour
$j$, the formula $\bigvee_{S\subseteq M}\neg x_{S,j}$, which says that the subsets of $M$ do not all have colour $j$.

A finite subset of $\Gamma$ mentions only finitely many integers, all at most some $N$. A bad colouring $c$ of $\{1,\dots,N\}$ satisfies it, setting $x_{S,j}$ true when
$c(S)=j$ (and the variables of sets not inside $\{1,\dots,N\}$ as we like). By the compactness theorem (Theorem~\ref{thm:compactness}), $\Gamma$ has a model $\alpha$. Colour
each $k$-subset $S$ of $\N^*$ with the least $j$ such that $\alpha(x_{S,j})=1$. By Theorem~\ref{thm:ramseyinf} there is an infinite homogeneous set, say of colour $j$, and any
$m$ of its elements form a set $M$ all of whose $k$-subsets $S$ have $\alpha(x_{S,j})=1$: the formula $\bigvee_{S\subseteq M}\neg x_{S,j}$ of $\Gamma$ is false, a contradiction.
\end{proof}

This proof gives no bound at all on $N$: compactness only says that a finite $N$ exists. A direct proof (Exercise~\ref{ex:ramseydirect}) gives a tower of
exponentials, and the truth is close to that for $k\ge3$. In 1977 Jeff Paris and Leo Harrington found a stunning variant: require in addition that the homogeneous
set $H$ be \emph{large}, $|H|\ge\min H$. The finite statement is still true, by the same compactness argument. But it cannot be proved in Peano arithmetic, the
usual axioms for the natural numbers: the least $N$ grows faster than any function the axioms can prove total. It was the first natural mathematical
statement shown to be independent of arithmetic, in the sense of G\"odel, and Tome~C3 explains what such a sentence means.

Ramsey's theorem for $4$-subsets proves the result of Erd\H os and Szekeres that opened Chapter~2 (Exercise~\ref{ex:happyending}): among enough points in general
position, some $n$ form a convex polygon. And for pairs, it contains Theorem~\ref{thm:es} in a qualitative form (Exercise~\ref{ex:monoseq}).

\section{Colouring the integers}

In 1916, studying Fermat's last theorem modulo primes, Issai Schur found one of the first Ramsey-type theorems, twelve years before Ramsey.

\begin{theorem}[Schur, 1916]\label{thm:schur}
Let $r\ge1$ and $N\ge R_r(3)-1$. For every colouring of $\{1,\dots,N\}$ with $r$ colours there are $x,y,z$ of the same colour with $x+y=z$ ($x=y$ allowed).
\end{theorem}
\begin{proof}
Colour the edge $\{i,j\}$ of the complete graph on $\{0,1,\dots,N\}$ with the colour of $|i-j|$. This complete graph has $N+1\ge R_r(3)$ vertices, so it has a monochromatic
triangle $i<j<k$. Then $x=j-i$, $y=k-j$ and $z=k-i$ have the same colour and $x+y=z$.
\end{proof}

The least $N$ for which the conclusion holds, minus one, is the \emph{Schur number} $S(r)$: the largest $N$ such that $\{1,\dots,N\}$ can be split into $r$ sum-free
sets. The values $S(1)=1$, $S(2)=4$, $S(3)=13$, $S(4)=44$ are classical; $S(5)=160$ was established by Marijn Heule in 2017 with a SAT solver, in a proof of two petabytes,
the largest proof ever written at the time. SAT solvers, Chapter~1's subject, are now the main tool for exact Ramsey-type numbers.

\begin{corollary}[Dickson, 1909; Schur, 1916]\label{cor:schurfermat}
Let $m\ge1$. For every prime $p>R_m(3)$, the congruence $x^m+y^m\equiv z^m\pmod p$ has a solution with $x,y,z$ not divisible by $p$.
\end{corollary}
\begin{proof}
In the group $\mathbb F_p^*$, the $m$-th powers form the image $H$ of the homomorphism $u\mapsto u^m$. Its kernel is the set of roots of $u^m-1$, at most $m$ elements in the field
$\mathbb F_p$, so $H$ has index at most $m$: there are $r\le m$ cosets $a_1H,\dots,a_rH$. Colour each $n\in\{1,\dots,p-1\}$ by the coset of $n\bmod p$. Since
$p-1\ge R_m(3)\ge R_r(3)-1$, Schur's theorem gives $x+y=z$ in $\{1,\dots,p-1\}$ with $x,y,z$ in the same coset $aH$: $x\equiv au^m$, $y\equiv av^m$, $z\equiv aw^m$. Dividing by $a$,
$u^m+v^m\equiv w^m\pmod p$, with $u,v,w$ invertible.
\end{proof}

So Fermat's equation has non-trivial solutions modulo every large prime, and no argument that works modulo primes alone can prove Fermat's last theorem.
Leonard Dickson had proved this in 1909 by long computations with cyclotomy; Schur's proof is the one above. The truth came, in 1994, from elliptic curves and modular forms (Tome~12 tells part of the story).

Now arithmetic progressions. A \emph{$k$-term arithmetic progression} (or $k$-AP) is a set $\{a,a+d,\dots,a+(k-1)d\}$ with $d\ge1$.

\begin{theorem}[van der Waerden, 1927]\label{thm:vdw}
For all $k,r\ge1$ there is an integer $W(k,r)$ such that every colouring of $\{1,\dots,W(k,r)\}$ with $r$ colours contains a monochromatic $k$-AP.
\end{theorem}

The conjecture came from Baudet and Schur; van der Waerden heard it in Hamburg and proved it in a few days, in conversations with Artin and Schreier that
he later described in a famous essay on how mathematical ideas are found. The proof below is his, reorganised as \emph{colour focusing}. The idea is to build,
not one progression, but a bouquet of progressions of length $k-1$ of different colours all aiming at the same next point, their \emph{focus}. When the bouquet has
one progression of each colour, the focus, whatever its colour, completes one of them.

Write $[N]=\{1,\dots,N\}$. Monochromatic $(k-1)$-APs $\{a_i+jd_i:0\le j\le k-2\}$, $i=1,\dots,s$, are \emph{colour-focused at $f$} if they have distinct colours and $a_i+(k-1)d_i=f$
for all $i$.

\begin{lemma}[colour focusing]\label{lem:focus}
Let $k\ge2$, $r\ge1$, and suppose $W(k-1,r')$ exists for every $r'$. For each $s\in\{1,\dots,r\}$ there is an $N(s)$ such that every colouring of $[N(s)]$ with $r$ colours
contains a monochromatic $k$-AP, or $s$ colour-focused $(k-1)$-APs whose focus lies in $[N(s)]$.
\end{lemma}
\begin{proof}
Induction on $s$. For $s=1$ take $N(1)=2W(k-1,r)$: the first half $[W(k-1,r)]$ contains a monochromatic $(k-1)$-AP $a,\dots,a+(k-2)d$, and its focus satisfies
$a+(k-1)d\le W(k-1,r)+d\le2W(k-1,r)$.

Let $s<r$, let $N=N(s)$, and set
\[
N(s+1)=N\cdot2W\bigl(k-1,r^{N}\bigr).
\]
Cut $[N(s+1)]$ into consecutive \emph{blocks} $B_1,B_2,\dots$ of length $N$, $B_j=\{N(j-1)+1,\dots,Nj\}$. The colouring of a block, read from left to right, is one of $r^{N}$
patterns; colour each block index $j$ with the pattern of $B_j$. Among the first $W(k-1,r^{N})$ indices there is a monochromatic $(k-1)$-AP $b,b+e,\dots,b+(k-2)e$: the
blocks $B_b,B_{b+e},\dots,B_{b+(k-2)e}$ are coloured identically, and since $b+(k-2)e\le W(k-1,r^{N})$ and $e\le W(k-1,r^{N})$, also $b+(k-1)e\le2W(k-1,r^{N})$, so the block
$B_{b+(k-1)e}$ exists.

Apply the induction hypothesis to $B_b$, a translate of $[N]$. If it contains a monochromatic $k$-AP, we are done. Otherwise it contains $(k-1)$-APs
$A_i=\{a_i+jd_i\}$, $i=1,\dots,s$, of distinct colours $c_1,\dots,c_s$, colour-focused at a point $f\in B_b$. If $f$ has one of the colours $c_i$, then $A_i\cup\{f\}$ is a monochromatic $k$-AP. Otherwise $f$ has a new colour $c_{s+1}$. Let $D=Ne$ be the distance between
consecutive chosen blocks, and consider
\[
A'_i=\{a_i+j(d_i+D):0\le j\le k-2\}\ (i\le s),\qquad A'_{s+1}=\{f+jD:0\le j\le k-2\}.
\]
The $j$-th term of $A'_i$ is the $j$-th term of $A_i$ moved $j$ blocks of the progression to the right, into $B_{b+je}$, at the same position in the block; since these blocks have
the same pattern, it has colour $c_i$. Likewise $A'_{s+1}$ consists of the copies of $f$ in $B_b,\dots,B_{b+(k-2)e}$, all of colour $c_{s+1}$. Each $A'_i$ has the same next term
\[
a_i+(k-1)(d_i+D)=f+(k-1)D=f',
\]
which lies in $B_{b+(k-1)e}\subseteq[N(s+1)]$. These are $s+1$ colour-focused $(k-1)$-APs.
\end{proof}

\begin{proof}[Proof of Theorem~\ref{thm:vdw}]
Induction on $k$. For $k=1$, $W(1,r)=1$; for $k=2$, $W(2,r)=r+1$ by the pigeonhole principle. For $k\ge3$ apply the lemma with $s=r$: in $[N(r)]$, either there is a
monochromatic $k$-AP, or there are $r$ colour-focused $(k-1)$-APs, one of each colour, and the focus has the colour of one of them, which it completes into a
monochromatic $k$-AP. So $W(k,r)\le N(r)$.
\end{proof}

\begin{figure}[ht]
\centering
\fig{ch08-focusing}
\caption{Colour focusing for $k=3$. Top: inside one block, a red and a blue $2$-AP both aim at the same point $f$, which is green. Bottom: in identically
coloured blocks $B_b$ and $B_{b+e}$, the red and blue progressions are stretched by the block distance $D$, and the copies of $f$ give a green one; all three now aim
at $f'=f+2D$.}
\label{fig:focusing}
\end{figure}

The bound this proof gives grows like Ackermann's function of Chapter~2: $W(k,2)$ is bounded by a tower of towers. In 1988 Saharon Shelah found a proof with
primitive recursive bounds, and in 2001 Timothy Gowers proved
\[
W(k,r)\le2^{2^{r^{2^{2^{k+9}}}}}
\]
as a by-product of his Fourier-analytic proof of Szemer\'edi's theorem. The exact values are rare.

\begin{proposition}\label{prop:w32}
$W(3,2)=9$.
\end{proposition}
\begin{proof}
The colouring $RRBBRRBB$ of $[8]$ has no monochromatic $3$-AP: the red numbers are $1,2,5,6$ and the blue ones $3,4,7,8$, and in each set the middle of two elements, when
it is an integer, is not in the set ($\frac{1+5}2=3$, $\frac{2+6}2=4$, and the other sums are odd; the blue set is the red one shifted by $2$).

Now colour $[9]$ with no monochromatic $3$-AP, and say $5$ is red. For each $x$, the numbers $x$, $5$, $10-x$ form a $3$-AP, so $x$ and $10-x$ are not both red. If $4$ and $6$ are
both blue, then $2$ (from $2,4,6$) and $8$ (from $4,6,8$) are red, and $2,5,8$ is a red $3$-AP. Otherwise exactly one of $4,6$ is red; by the symmetry $x\mapsto10-x$, say $4$
is red and $6$ blue. Then, one forced colour after another: $3$ is blue ($3,4,5$); $9$ is red ($3,6,9$); $1$ is blue ($1,5,9$); $7$ is blue ($5,7,9$); $8$ is red ($6,7,8$);
$2$ is blue ($2,5,8$). Now $1,2,3$ is a blue $3$-AP.
\end{proof}

The other known values are $W(3,3)=27$, $W(4,2)=35$, $W(3,4)=76$, $W(5,2)=178$ and $W(6,2)=1132$ (Kouril and Paul, 2008, with a SAT solver on FPGA hardware).

\begin{remark}
Van der Waerden's theorem says that one of the colour classes contains long progressions. Erd\H os and Tur\'an asked in 1936 whether \emph{every} set of positive upper density
does. Szemer\'edi proved it in 1975: for every $\delta>0$ and $k$, every subset of $[N]$ of size $\delta N$ contains a $k$-AP when $N$ is large enough. Four completely different
proofs followed, by ergodic theory (Furstenberg, 1977), Fourier analysis (Gowers), hypergraph regularity, and the density Hales--Jewett theorem (Polymath, 2009). The primes
have density $0$, but Ben Green and Terence Tao proved in 2004 that they contain arbitrarily long progressions, such as the $26$ primes
$43\,142\,746\,595\,714\,191+23\,681\,770\cdot23\#\cdot n$, $n=0,\dots,25$, found in 2010.
\end{remark}

\section{Sunflowers}

A \emph{sunflower} with $p$ petals is a family of $p$ distinct sets $A_1,\dots,A_p$ whose pairwise intersections are all equal to the same set $Y$, the \emph{core}. The
\emph{petals} $A_i\setminus Y$ are then pairwise disjoint. Pairwise disjoint sets form a sunflower with empty core.

\begin{theorem}[Erd\H os and Rado, 1960]\label{thm:sunflower}
Let $k\ge1$ and $p\ge2$. A family of more than $k!(p-1)^k$ sets, each of size at most $k$, contains a sunflower with $p$ petals.
\end{theorem}
\begin{proof}
We first prove it for sets of size exactly $k$, by induction on $k$. For $k=1$, the family consists of more than $p-1$ distinct singletons, which are pairwise disjoint: any $p$ of them form a sunflower.

Let $k\ge2$ and $\mathcal F$ a family of more than $k!(p-1)^k$ sets of size $k$. Choose a maximal subfamily $\mathcal D\subseteq\mathcal F$ of pairwise disjoint sets. If $|\mathcal D|\ge p$,
any $p$ sets of $\mathcal D$ form a sunflower. Otherwise the union $U$ of $\mathcal D$ has at most $k(p-1)$ elements, and by maximality every set of $\mathcal F$ meets $U$. So some $x\in U$ belongs to
more than
\[
\frac{k!(p-1)^k}{k(p-1)}=(k-1)!(p-1)^{k-1}
\]
sets of $\mathcal F$. Remove $x$ from the sets containing it: they become more than $(k-1)!(p-1)^{k-1}$ distinct sets of size $k-1$. By the induction hypothesis, $p$ of them form a sunflower
$A_1\setminus\{x\},\dots,A_p\setminus\{x\}$ with core $Y$; then $A_1,\dots,A_p$ form a sunflower with core $Y\cup\{x\}$.

For sets of size at most $k$, add to each set $A$ new elements, different for different sets, until it has size $k$. The new family has the same number of sets and
contains a sunflower with $p$ petals; the added elements belong to only one set each, so they are not in the core, and removing them leaves $p$ distinct sets whose pairwise
intersections are still the core.
\end{proof}

The bound is not far from the truth: there are families of $(p-1)^k$ sets of size $k$ with no sunflower with $p$ petals (Exercise~\ref{ex:sunflowersharp}). Erd\H os
conjectured that $C(p)^k$ sets always suffice, and offered \$1000 for a proof. In 2019 Alweiss, Lovett, Wu and Zhang brought the bound down to about $(Cp\log k)^k$, an
improvement that came from an unexpected link with random restrictions in complexity theory; the conjecture is still open. The link runs both ways: in 1985 Alexander
Razborov used sunflowers to prove that detecting a clique of size $k$ in a graph needs monotone Boolean circuits of super-polynomial size. Tome~C3 tells that story.

\begin{beauty}
The theorems of this chapter all share one shape: a large enough object, \emph{however} it is cut, has a large regular piece. Ramsey cut pairs, van der Waerden cut
integers, Erd\H os and Rado cut families of sets. The proofs share a shape too: find a point that sees many others in the same way, restrict to them, and repeat;
or move up one level, and colour blocks by their patterns. And the numbers that measure the phenomenon are, with very few exceptions, out of reach: we can prove that
order must appear long before we can say exactly when.
\end{beauty}

\section{Exercises}

\begin{exercise}[$\star\star$]\label{ex:turaninduction}
Give another proof of the bound of Theorem~\ref{thm:turan} by induction on $n$: a $K_{r+1}$-free graph with the most edges contains a $K_r$ (why?); count the edges inside it,
between it and the rest, and in the rest, and check that $t(n,r)=t(n-r,r)+(n-r)(r-1)+\binom r2$.
\end{exercise}

\begin{exercise}[$\star\star$]\label{ex:zarankiewicz}
\emph{Zarankiewicz's problem.} Show that an $n\times n$ matrix of zeros and ones in which no two rows have ones in two common columns has at most $\frac n2\bigl(1+\sqrt{4n-3}\bigr)$
ones.
\end{exercise}

\begin{exercise}[$\star$]\label{ex:polarity}
Draw the graph $G_2$ of Theorem~\ref{thm:polarity} (seven vertices) and count its edges. It is known that $\mathrm{ex}(7,C_4)=9$.
\end{exercise}

\begin{exercise}[$\star$]\label{ex:schurbound}
Deduce from Proposition~\ref{prop:r333} that $R_r(3)\le\lfloor e\,r!\rfloor+1$ for all $r\ge1$. Check the cases $r=2,3$.
\end{exercise}

\begin{exercise}[$\star\star\star$]\label{ex:r333lower}
\emph{Greenwood and Gleason's colouring.} Let $\mathbb F_{16}$ be the field with $16$ elements (Tome~2, Chapter~9). Its multiplicative group is cyclic of order $15$; let $H$ be the subgroup of
cubes, of order $5$, with cosets $H$, $gH$, $g^2H$. Colour the edge $\{x,y\}$ of the complete graph on $\mathbb F_{16}$ by the coset of $x-y$. Show that there is no monochromatic
triangle, so $R_3(3)=17$.
\end{exercise}

\begin{exercise}[$\star$]\label{ex:schur2}
Show that $S(2)=4$: split $\{1,2,3,4\}$ into two sum-free sets, and show that $\{1,\dots,5\}$ cannot be split.
\end{exercise}

\begin{exercise}[$\star$]\label{ex:fermatp}
Show that $x^3+y^3\equiv z^3\pmod{13}$ has no solution with $x,y,z$ prime to $13$, and relate this to the colouring of $K_{13}$ in Theorem~\ref{thm:smallramsey}. (The notebook finds all such
primes: $7$ and $13$ are the only ones.)
\end{exercise}

\begin{exercise}[$\star$]\label{ex:monoseq}
Using Theorem~\ref{thm:ramseyinf} for pairs, show that every sequence of real numbers has an infinite monotone subsequence. Compare with Theorem~\ref{thm:es}.
\end{exercise}

\begin{exercise}[$\star\star$]\label{ex:happyending}
\emph{The happy ending problem.} (a) Show that finitely many points in the plane, no three on a line, are the vertices of a convex polygon as soon as any four of them are
the vertices of a convex quadrilateral. (b) Using Exercise~\ref{ex:klein} and Theorem~\ref{thm:ramseyfin} for $4$-subsets, show that for every $n$ there is an $N$ such that among any
$N$ points in general position, some $n$ are the vertices of a convex $n$-gon.
\end{exercise}

\begin{exercise}[$\star\star$]\label{ex:ramseydirect}
Prove Theorem~\ref{thm:ramseyfin} directly, by copying the proof of Theorem~\ref{thm:ramseyinf} with finite sets, and get an explicit bound for $R^{(k+1)}_r(m)$ in terms of the
function $n\mapsto R^{(k)}_r(n)$. How fast does your bound grow?
\end{exercise}

\begin{exercise}[$\star$]\label{ex:wlower}
Using a random colouring as in Tome~3, Theorem~12.2, show that $W(k,2)>\sqrt{k-1}\;2^{k/2}$ for $k\ge2$. (Count the $k$-APs in $[N]$ first.)
\end{exercise}

\begin{exercise}[$\star$]\label{ex:sunflowersharp}
For $k\ge1$ and $p\ge2$, consider the $(p-1)^k$ sets $\{(1,a_1),(2,a_2),\dots,(k,a_k)\}$ with $a_1,\dots,a_k\in\{1,\dots,p-1\}$. Show that no $p$ of them form a sunflower.
\end{exercise}

\section{Hints}

\begin{hint}{ex:turaninduction}
Adding an edge to a graph without $K_r$ cannot create a $K_{r+1}$. A vertex outside a $K_r$ is adjacent to at most $r-1$ of its vertices.
\end{hint}
\begin{hint}{ex:zarankiewicz}
Count pairs of ones in the same row, as in Theorem~\ref{thm:kst} with $s=t=2$, but keep $\binom n2$ exactly.
\end{hint}
\begin{hint}{ex:polarity}
Over $\mathbb F_2$, a vector is orthogonal to itself when it has an even number of ones.
\end{hint}
\begin{hint}{ex:schurbound}
Set $u_r=R_r(3)-1$ and divide the recursion by $r!$.
\end{hint}
\begin{hint}{ex:r333lower}
A monochromatic triangle gives $a,b\in H$ with $a\ne b$ and $a+b\in H$ after a division. Show $(1+a)^5\ne1$ for $a\in H$, $a\ne1$, using $a^4+a^3+a^2+a+1=0$ and $(1+a)^4=1+a^4$.
\end{hint}
\begin{hint}{ex:schur2}
If $1$ is red, $2$ must be blue.
\end{hint}
\begin{hint}{ex:fermatp}
List the non-zero cubes modulo $13$.
\end{hint}
\begin{hint}{ex:monoseq}
Colour a pair $i<j$ by whether $x_i<x_j$.
\end{hint}
\begin{hint}{ex:happyending}
(a) A point inside the convex hull of others is inside a triangle formed by three of them. (b) Colour each $4$-set by convex or not.
\end{hint}
\begin{hint}{ex:ramseydirect}
Build $x_1,x_2,\dots$ and nested sets as in the infinite proof, keeping track of sizes; you need $m$ indices with the same $\chi_i$.
\end{hint}
\begin{hint}{ex:wlower}
A $k$-AP in $[N]$ is determined by its first term and its difference $d\le\frac{N-1}{k-1}$.
\end{hint}
\begin{hint}{ex:sunflowersharp}
In a sunflower, what can happen at a fixed coordinate $i$?
\end{hint}

\section{Solutions}

\begin{solution}{ex:turaninduction}
\idea Remove a whole $K_r$ at once: the rest is a smaller instance, and the edges between are controlled by $K_{r+1}$-freeness.
\writeup Induction on $n$; for $n\le r$, $t(n,r)=\binom n2$ and there is nothing to prove. Let $n>r$ and $G$ be $K_{r+1}$-free with the most edges. If $G$ had no $K_r$, it would not be
complete, and adding a missing edge $uv$ would create no $K_{r+1}$: a clique containing $u$ and $v$ minus $u$ is a clique of $G$, of size at most $r-1$. That contradicts maximality, so $G$
contains a clique $K$ with $r$ vertices. Each vertex outside $K$ is adjacent to at most
$r-1$ vertices of $K$, otherwise it would form a $K_{r+1}$ with them. The graph $G-K$ has no $K_{r+1}$, so by induction it has at most $t(n-r,r)$ edges (for $n-r<r$, at most
$\binom{n-r}2=t(n-r,r)$). So $e(G)\le\binom r2+(n-r)(r-1)+t(n-r,r)$. Removing one vertex from each class of $T(n,r)$ gives $T(n-r,r)$; the removed vertices form a $K_r$,
and each remaining vertex is adjacent to $r-1$ of them, which proves the identity, and $e(G)\le t(n,r)$.
\end{solution}

\begin{solution}{ex:zarankiewicz}
\idea Two rows share at most one pair of columns, so the pairs of ones in the rows are all different.
\writeup Let row $i$ have $r_i$ ones, $m=\sum r_i$. Each row contains $\binom{r_i}2$ pairs of columns where it has ones, and no pair of columns appears in two rows, so
$\sum_i r_i(r_i-1)\le n(n-1)$. By Cauchy--Schwarz, $\sum r_i^2\ge\frac{m^2}n$, hence $\frac{m^2}n-m\le n(n-1)$, that is $m^2-nm-n^2(n-1)\le0$. The positive root of this quadratic is
$\frac12\bigl(n+\sqrt{n^2+4n^2(n-1)}\bigr)=\frac n2\bigl(1+\sqrt{4n-3}\bigr)$. The incidence matrix of points and lines of a projective plane of order $q$ ($n=q^2+q+1$, $q+1$ ones per row) attains the bound.
\end{solution}

\begin{solution}{ex:polarity}
\idea List the seven non-zero vectors of $\mathbb F_2^3$ and test the dot products.
\writeup Each point $x$ has the $3$ points of $x^\perp$ as neighbours, minus $x$ itself when $x\cdot x=0$, which over $\mathbb F_2$ means $x$ has an even number of ones: $110$, $101$, $011$. These
three have degree $2$, the four others ($100$, $010$, $001$, $111$) degree $3$. The number of edges is $\frac12(3\cdot2+4\cdot3)=9=\mathrm{ex}(7,C_4)$: for instance $100$ is adjacent to $010$, $001$, $011$, and
$111$ to $110$, $101$, $011$. The construction is optimal here.
\end{solution}

\begin{solution}{ex:schurbound}
\idea The recursion $u_r\le ru_{r-1}+1$ becomes a telescoping sum after division by $r!$, and the sum is a partial sum of the series of $e$.
\writeup With $u_r=R_r(3)-1$, Proposition~\ref{prop:r333} reads $u_r\le ru_{r-1}+1$, and $u_1=2$ (one colour: any $3$ vertices). Divide by $r!$: $\frac{u_r}{r!}\le\frac{u_{r-1}}{(r-1)!}+\frac1{r!}$, so
$\frac{u_r}{r!}\le2+\frac1{2!}+\dots+\frac1{r!}=\sum_{j=0}^r\frac1{j!}$. The number $r!\sum_{j=0}^r\frac1{j!}$ is an integer, and $e\,r!$ exceeds it by $r!\sum_{j>r}\frac1{j!}<\frac1{r+1}+\frac1{(r+1)^2}+\cdots=\frac1r\le1$, so it equals
$\lfloor e\,r!\rfloor$. Hence $R_r(3)\le\lfloor e\,r!\rfloor+1$: $6$ for $r=2$ and $17$ for $r=3$, both exact.
\end{solution}

\begin{solution}{ex:r333lower}
\idea In characteristic $2$, $x-y=y-x=x+y$, so the colour is well defined, and a triangle $x,y,z$ has $(x+z)=(x+y)+(y+z)$: a monochromatic triangle is a solution of $u+v=w$ in one coset.
\writeup The cubes form the image of $u\mapsto u^3$ in the cyclic group of order $15$, a subgroup of order $5$ and index $3$, and $H=\{y:y^5=1\}$. A monochromatic triangle $x,y,z$ gives $u=x+y$ and
$v=y+z$ in the same coset $cH$, with $u+v=x+z$ also in $cH$ and $u\ne v$ (as $x\ne z$). Divide by $u$: $a=\frac vu\in H$, $a\ne1$, and $1+a=\frac{u+v}u\in H$. Now $a$ has order $5$, so
$a^4+a^3+a^2+a+1=\frac{a^5-1}{a-1}=0$. In characteristic $2$, $(1+a)^2=1+a^2$ and $(1+a)^4=1+a^4$, so $(1+a)^5=(1+a^4)(1+a)=1+a+a^4+a^5=a+a^4$. If this were $1$, then $a^4+a+1=0$; subtracting from
$a^4+a^3+a^2+a+1=0$ gives $a^3+a^2=a^2(a+1)=0$, so $a=0$ or $a=1$, both excluded. So $(1+a)^5\ne1$, $1+a\notin H$: no monochromatic triangle, and $R_3(3)>16$. With Proposition~\ref{prop:r333},
$R_3(3)=17$. Each colour class is a $5$-regular graph on $16$ vertices without triangles, the \emph{Clebsch graph}.
\end{solution}

\begin{solution}{ex:schur2}
\idea Start with $1$ and follow the forced choices.
\writeup $\{1,4\}$ and $\{2,3\}$ are sum-free: $1+1=2$, $1+4=5$, $4+4=8$ are not in the first, $2+2=4$, $2+3=5$, $3+3=6$ not in the second. For $\{1,\dots,5\}$, say $1$ is red. Then $2$ is blue
($1+1=2$), so $4$ is red ($2+2=4$), so $3$ is blue ($1+3=4$). Now $5=1+4$ forbids red and $5=2+3$ forbids blue.
\end{solution}

\begin{solution}{ex:fermatp}
\idea The colouring by cubic residues mod $13$ is exactly the colouring with no red triangle used for $R(3,5)$.
\writeup The cubes of $1,\dots,12$ modulo $13$ are $1,8,1,12$, $8,8,5,5$, $1,12,5,12$, so the non-zero cubes form $H=\{1,5,8,12\}=\{\pm1,\pm5\}$. A solution with $xyz$ prime to $13$ gives $a+b\equiv c$ with $a,b,c\in H$. But
the proof of Theorem~\ref{thm:smallramsey} showed that sums of two elements of $\{1,5,8,12\}$ are never in the set. So the red graph of the colouring of $K_{13}$ used there is
the colour class of the cubes in Schur's colouring by cosets: Schur's colouring of $\{1,\dots,12\}$ has no monochromatic solution, and indeed Corollary~\ref{cor:schurfermat} only promises
one for $p>R_3(3)=17$. For $p=7$ the non-zero cubes are $\{1,6\}=\{\pm1\}$, and $1+1$, $1+6$, $6+6$ are $2,0,5$: no solution either.
\end{solution}

\begin{solution}{ex:monoseq}
\idea Two colours on pairs: rising or not. An infinite homogeneous set is a monotone subsequence.
\writeup Colour $\{i,j\}$, $i<j$, red if $x_i<x_j$ and blue otherwise. An infinite homogeneous set $i_1<i_2<\cdots$ gives a strictly increasing subsequence (red) or a non-increasing one (blue).
Theorem~\ref{thm:es} is the quantitative finite version: $N>(a-1)(b-1)$ distinct terms contain an increasing subsequence of length $a$ or a decreasing one of length $b$, while Ramsey's theorem for pairs only gives
$R(a,b)$, which is far larger.
\end{solution}

\begin{solution}{ex:happyending}
\idea Being in convex position is a property of $4$-sets; Klein's lemma rules out a large homogeneous set of the colour ``not convex''.
\writeup (a) If the points are not in convex position, one of them, $q$, lies in the convex hull of the others. Triangulate that convex polygon from one of its vertices: $q$ lies in one of the triangles,
and as no three points are collinear, strictly inside; then $q$ and the three vertices of that triangle are four points not in convex position.
(b) Colour each $4$-subset of the $N$ points red if it is in convex position, blue otherwise, with $N=R^{(4)}_2(\max(n,5))$. A blue homogeneous set of $5$ points would contradict Exercise~\ref{ex:klein}, so there is a red homogeneous set of $n$
points, and by (a) they form a convex $n$-gon. Erd\H os and Szekeres's 1935 proof was exactly this, with Ramsey's theorem rediscovered by Szekeres.
\end{solution}

\begin{solution}{ex:ramseydirect}
\idea In the finite version, each step of the construction keeps a homogeneous part of the remaining points, and we need about $rm$ steps.
\writeup Let $f(n)=R^{(k)}_r(n)$, with $f(n)=n$ for $n<k$. We build $L=r(m-1)+2$ points $x_1,\dots,x_L$. Define $n_{L-1}=1$ and
$n_{i-1}=f(n_i)+1$ for $i=L-1,\dots,1$, and let $N=n_0$. Starting from $X_0=[N]$, at step $i\le L-1$ choose $x_i\in X_{i-1}$; the set $X_{i-1}\setminus\{x_i\}$ has $f(n_i)$ elements, so it contains a set $X_i$ of
size $n_i$ on which the induced colouring of $k$-subsets is constant, of colour $\chi_i$. Finally $x_L$ is the element of $X_{L-1}$. As in the infinite proof, a $(k+1)$-subset of $\{x_1,\dots,x_L\}$ has the colour
$\chi_i$ of its element of smallest index, so $m$ points $x_i$ with $i<L$ and equal $\chi_i$ form a homogeneous set, and they exist by the pigeonhole principle, as there are $L-1=r(m-1)+1$ colours $\chi_1,\dots,\chi_{L-1}$. Then $R^{(k+1)}_r(m)\le n_0$, the composition of $n\mapsto f(n)+1$ about $rm$ times. For pairs ($k=1$, $f(n)=r(n-1)+1$) this is exponential, of order $r^{rm}$; for triples a
tower of two exponentials, and each new $k$ adds a level. The stepping-up lemma of Erd\H os, Hajnal and Rado shows that towers of exponentials really occur.
\end{solution}

\begin{solution}{ex:wlower}
\idea The first moment method: if the expected number of monochromatic $k$-APs is less than $1$, some colouring has none.
\writeup A $k$-AP in $[N]$ is given by $a\ge1$ and $d\ge1$ with $a+(k-1)d\le N$; for each $d$ there are $N-(k-1)d$ choices of $a$, so the number is $\sum_{d\ge1}(N-(k-1)d)^+<\frac{N^2}{2(k-1)}$. In a uniform random
$2$-colouring, each is monochromatic with probability $2^{1-k}$. If $N^2\le(k-1)2^k$, the expectation is below $\frac{N^2}{2(k-1)}2^{1-k}\le1$, and some colouring of $[N]$ has no monochromatic $k$-AP:
$W(k,2)>N$ for $N=\lfloor\sqrt{k-1}\,2^{k/2}\rfloor$. Berlekamp (1968) did better with finite fields, $W(p+1,2)>p2^p$ for prime $p$.
\end{solution}

\begin{solution}{ex:sunflowersharp}
\idea Coordinate by coordinate, a sunflower either agrees everywhere or is pairwise different; with only $p-1$ values, it must agree.
\writeup Suppose $A_1,\dots,A_p$ is a sunflower with core $Y$, and fix a coordinate $i$. If two of the sets share the element $(i,a)$, it lies in their intersection, which is $Y$, so it lies in all $p$ sets. So at
coordinate $i$, the values $a_i$ of the $p$ sets are either all equal or pairwise distinct. Pairwise distinct is impossible with only $p-1$ values. So all the sets agree at every coordinate: they are equal,
which contradicts distinctness.
\end{solution}

\section*{Sources and further reading}

\begin{itemize}
\item F.~P.~Ramsey, \emph{On a problem of formal logic}, Proceedings of the London Mathematical Society 30 (1930), 264--286; I.~Schur, \emph{\"Uber die Kongruenz
$x^m+y^m\equiv z^m\pmod p$}, Jahresbericht der DMV 25 (1916), 114--117; B.~L.~van der Waerden, \emph{Beweis einer Baudetschen Vermutung}, Nieuw Archief voor
Wiskunde 15 (1927), 212--216, and \emph{How the proof of Baudet's conjecture was found}, in Studies in Pure Mathematics, Academic Press, 1971, 251--260.
\item R.~L.~Graham, B.~L.~Rothschild and J.~H.~Spencer, \emph{Ramsey Theory}, 2nd ed., Wiley, 1990: the standard reference, with colour focusing; A.~Soifer,
\emph{The Mathematical Coloring Book}, Springer, 2009, for the history of Schur, van der Waerden and Ramsey.
\item P.~Tur\'an, \emph{On an extremal problem in graph theory} (in Hungarian), Matematikai \'es Fizikai Lapok 48 (1941), 436--452; A.~A.~Zykov, \emph{On some properties of
linear complexes}, Matematicheskii Sbornik 24 (1949), 163--188; M.~Aigner and G.~M.~Ziegler, \emph{Proofs from THE BOOK}, 6th ed., Springer, 2018, chapter on Tur\'an's
theorem, with five proofs.
\item T.~K\H ov\'ari, V.~T.~S\'os and P.~Tur\'an, \emph{On a problem of K.~Zarankiewicz}, Colloquium Mathematicum 3 (1954), 50--57; W.~G.~Brown, \emph{On graphs that do not contain
a Thomsen graph}, Canadian Mathematical Bulletin 9 (1966), 281--285; Z.~F\"uredi and M.~Simonovits, \emph{The history of degenerate (bipartite) extremal graph problems},
in Erd\H os Centennial, Springer, 2013, 169--264.
\item R.~E.~Greenwood and A.~M.~Gleason, \emph{Combinatorial relations and chromatic graphs}, Canadian Journal of Mathematics 7 (1955), 1--7; S.~P.~Radziszowski, \emph{Small
Ramsey numbers}, Electronic Journal of Combinatorics, Dynamic Survey DS1 (regularly revised); V.~Angeltveit and B.~D.~McKay, \emph{$R(5,5)\le46$}, preprint, 2024.
\item J.~Paris and L.~Harrington, \emph{A mathematical incompleteness in Peano arithmetic}, in Handbook of Mathematical Logic, North-Holland, 1977, 1133--1142.
\item P.~Erd\H os and R.~Rado, \emph{Intersection theorems for systems of sets}, Journal of the London Mathematical Society 35 (1960), 85--90; R.~Alweiss, S.~Lovett, K.~Wu and
J.~Zhang, \emph{Improved bounds for the sunflower lemma}, Annals of Mathematics 194 (2021), 795--815: a research paper whose key idea, a random ``spread'' argument, is
readable after this chapter and Tome~3.
\item M.~J.~H.~Heule, \emph{Schur number five}, Proceedings of AAAI 2018, 6598--6606; B.~Green and T.~Tao, \emph{The primes contain arbitrarily long arithmetic progressions},
Annals of Mathematics 167 (2008), 481--547.
\end{itemize}
